\documentclass{tietreport}

\coursecode{UCS503 - Software Engineering}
\projecttitle{SV Car Customs}
\projectsubtitle{An E-Commerce Platform for Vehicle-Compatible Car Accessories}
\students{%
  \texttt{1024170242} Vishu Verma\\
  \texttt{1024170241} Raghu Verma\\
  \texttt{1024170229} Pratham Sharma\\
  \texttt{1024170231} Tarush Mohindru}
\group{Group: \texttt{3Q24} BE Third Year, COPC}
\supervisor{Ms.\ Anushka}
\evaluation{Mid-Semester Evaluation (Prototype Stage) -- revised}

\begin{document}
\maketitle
\pagenumbering{roman}
\tableofcontents
\cleardoublepage
\pagenumbering{arabic}

\chapter{Introduction}

\paragraph{About this revision.}
This is the prototype-stage report submitted at the mid-semester evaluation,
re-checked against the code as it stood on 15 October 2026 (repository
revision \texttt{fb14d3d}, 75 commits). Statements that the first version
could only quote from notes are now verified in the code. The corrections are:
order states now use the real \texttt{order\_status} values; the GST rate is
an admin setting (default 28\%), not a fixed 18\%; free shipping applies from
the threshold \emph{inclusive}; project-scale figures are measured from the
repository; and the team member names are completed.

\section{Background and Context}

SV Car Customs is an e-commerce project for a Delhi business that sells
aftermarket car accessories, including body kits, spoilers, diffusers,
exhausts, lighting and steering wheels. Its central requirement is to connect
product discovery with vehicle compatibility, so a customer can find
accessories that fit a particular car before ordering. At the prototype stage
the project combined a customer storefront, an administration panel and a
separate \texttt{/customs} page presenting the business's vehicle
customisation services.

The work completed by mid-semester covered catalogue and SKU management,
vehicle fitment, customer authentication, checkout, PayU payments, order
administration, operational monitoring and configurable branding, deployed
as two independent applications (frontend and API).

\subsection{Problem Statement}

A generic product catalogue cannot express whether an accessory fits a
vehicle make and model. Customers also need accurate prices, reliable order
records and a payment flow that cannot be manipulated through
browser-submitted totals. Staff need one interface for maintaining products,
variants, stock and order status instead of disconnected records.

The project therefore builds a vehicle-aware commerce system in which the API
validates checkout data, each order keeps its historical details, and staff
manage the catalogue and fulfilment workflow. Damage-claim handling and
courier integration are part of the domain design but were not planned for
the prototype.

\section{Project Objectives}

\begin{enumerate}
  \item \textbf{Objective 1:} Catalogue browsing, product search and
        vehicle-fitment discovery through a responsive storefront.
  \item \textbf{Objective 2:} Customer authentication and separate
        administrative access, with server-side ownership and authorisation
        checks.
  \item \textbf{Objective 3:} SKU-based cart, trusted server-side price
        calculation and PayU-verified online payment, including Cash on
        Delivery (COD) with an online advance.
  \item \textbf{Objective 4:} Admin tools for products, categories, fitment,
        orders, home-page content and settings.
  \item \textbf{Objective 5:} Document the system through consistent UML
        views and establish functional, security and operational testing.
\end{enumerate}

These objectives defined the mid-semester scope. Quantified latency and
user-study targets were set later in the SRS and are evaluated in the final
report.

\chapter{Methodology and Design}

\section{System Architecture}

\subsection{High-Level Design}

The repository contains the Hono API in \texttt{svcar-api} and the Next.js
application in \texttt{svcar-india}. The two deploy independently and exchange
data through Zod-defined contracts shared in a \texttt{contracts} package.
The API runs on Fly.io, the frontend on Vercel, and Supabase provides
PostgreSQL and object storage (Figure~\ref{fig:arch}).

\begin{figure}[H]
\centering
\resizebox{\linewidth}{!}{%
\begin{tikzpicture}[node distance=9mm and 6mm]
  \node[ext] (cust) {Customer browser\\(storefront)};
  \node[ext, right=40mm of cust] (adm) {Owner / staff browser\\(admin panel)};
  \node[box, below=12mm of $(cust)!0.5!(adm)$, minimum width=72mm] (web) {\textbf{Vercel} -- Next.js 16 / React 19\\server-rendered, cached pages; session proxy};
  \node[box, below=10mm of web, minimum width=72mm] (api) {\textbf{Fly.io} -- Hono REST API (Node.js 20)\\JWT auth, Zod validation, pricing, orders};
  \node[ext, below left=12mm and -12mm of api] (db) {Supabase\\PostgreSQL + Storage};
  \node[ext, below=12mm of api] (payu) {PayU\\hosted checkout};
  \node[ext, below right=12mm and -12mm of api] (jobs) {Inngest, Resend,\\MSG91, Redis};
  \node[ext, left=8mm of api] (pin) {India Post\\pincode API};
  \node[ext, right=6mm of adm] (g) {Google Search /\\Merchant Center};
  \draw[arr] (cust) -- node[lbl]{HTTPS} (web);
  \draw[arr] (adm) -- node[lbl]{HTTPS} (web);
  \draw[arr] (web) -- node[lbl]{JSON / REST} (api);
  \draw[arr] (api) -- node[lbl]{Drizzle ORM} (db);
  \draw[arr,<->] (api) -- node[lbl]{hash + webhook} (payu);
  \draw[arr] (api) -- node[lbl]{events} (jobs);
  \draw[arr] (api) -- (pin);
  \draw[arr] (g) |- node[lbl,pos=0.75]{crawl, feed.xml} (web.east);
\end{tikzpicture}}
\caption{System architecture and external services}
\label{fig:arch}
\end{figure}

The API owns business validation and database access through Drizzle ORM.
PayU handles payment entry; its callbacks and webhooks return to the API for
verification. Inngest runs order-related background jobs, Resend sends
email, and Upstash Redis supports caching and rate limiting. Sentry and
structured logs provide operational visibility. The browser is outside the
trust boundary, and customer session routes are proxied through Next.js so
refresh cookies stay first-party.

\subsection{Technical Stack}

\begin{itemize}
  \item \textbf{Framework/Language:} TypeScript; Next.js 16, React 19 and
        Tailwind CSS 4 (frontend); Node.js 20 and Hono (backend).
  \item \textbf{Database:} Supabase PostgreSQL through Drizzle ORM; Upstash
        Redis for cache and rate limits.
  \item \textbf{Tools:} Git/GitHub, GitHub Actions, Vitest, Mock Service
        Worker, Docker, Vercel, Fly.io, Sentry.
  \item \textbf{Libraries:} Zustand (client state), SWR (fetching),
        react-hook-form and Zod (forms and validation), bcrypt, Inngest SDK.
\end{itemize}

\section{Implementation Details}

\subsection{Module 1: Core Functionality}

\paragraph{Catalogue and vehicle fitment.}
Products, images and categories form the catalogue; a junction table allows a
product in several categories. Dimensions, dimension values and SKUs model
saleable variants, so price and stock belong to the SKU. The fitment model
links products to vehicle makes and models; the ``My Garage'' selector
filters products by the shopper's car and product pages show fitment badges.

\paragraph{Authentication.}
Customers sign in with a six-digit email OTP, stored as a bcrypt hash, valid
for five minutes with at most five attempts. Verification issues a
short-lived access JWT and an httpOnly refresh cookie with a revocable
identifier; customers can later set a password. Admins use a separate
signing secret, and every admin request re-checks the admin flag in the
database (cached for 60 seconds). Password login compares against a dummy
bcrypt hash for unknown accounts to prevent timing-based account discovery.

\paragraph{Checkout and payment.}
At the prototype stage checkout required sign-in. The checkout endpoint
reloads the persisted cart, validates SKUs and calculates prices in integer
paise; the browser total is never used. It creates the order before payment,
then hands off to PayU with a signed hash. The callback and webhook verify
PayU's hash and the amount before confirming payment. An order-paid job
sends the confirmation email, and a reconcile job revisits orders stuck in
\emph{pending}. For COD, the customer pays a 15\% advance online and the
admin records collection of the balance.

\begin{table}[H]
\centering\small
\begin{tabularx}{\linewidth}{l X}
\toprule
\textbf{Quantity} & \textbf{Rule in the implementation} \\
\midrule
Subtotal & $S = \sum q_i \times p_i$ using current SKU prices from the database. \\
Shipping & $F = 0$ when $S \ge$ the free-shipping threshold (default \rupee5,000), otherwise the flat fee (default \rupee200). \\
GST & Prices are GST-inclusive: $G = S - S/(1+g)$ is shown, not added; $g$ is an admin setting (default 28\%). \\
Total & $T = S + F$. \\
COD & Advance $= 15\%$ of $T$ paid online; balance $= T -$ advance, collected on delivery. \\
\bottomrule
\end{tabularx}
\caption{Pricing rules}
\end{table}

\subsection{Module 2: Secondary Features}

\paragraph{Administration.}
The admin panel includes product editing with a variant matrix, fitment
editing and image uploads; category drag-and-drop with archive and restore;
orders with a status timeline and COD collection; a home-page builder and a
settings screen. Every admin mutation is written to an audit log. Further
views cover first-party analytics, WhatsApp leads, newsletter subscribers and
live activity.

\paragraph{Customs page and storefront delivery.}
The \texttt{/customs} page presented the customisation business with a
before/after slider, wrap configurator, builds carousel and booking call to
action. (It was later removed because it used placeholder content; see the
final report.) Home and product pages render on the server; a custom image
loader uses Supabase image rendering after the Vercel optimisation quota was
exhausted. SEO work included sitemap, robots, canonical URLs, structured data
and Open Graph images.

\paragraph{Brand configuration and tooling.}
One file at the repository root (\texttt{brand.config.ts}) defines the
brand. One command writes the derived configuration into both apps, and a
check command fails CI if the generated files drift. MSW fixtures let the frontend and admin run without a backend in
development.

\paragraph{Design scope not implemented.}
The domain model includes damage-claim submission and review and a courier
partner. These were kept as design elements only.

\subsection{Use Case Diagrams}
\begin{figure}[H]
\centering
\begin{tikzpicture}
  \draw[rounded corners=3pt] (-3.7,-6.4) rectangle (3.7,3.4);
  \node[font=\small\bfseries] at (0,3.05) {SV Car Customs};
  \node[uc] (u1) at (0,2.4) {Browse / search catalogue};
  \node[uc] (u2) at (0,1.6) {Check vehicle fitment};
  \node[uc] (u3) at (0,0.8) {Manage cart};
  \node[uc] (u4) at (0,0.0) {Checkout and pay};
  \node[uc] (u5) at (0,-0.8) {Sign in with email OTP};
  \node[uc] (u6) at (0,-1.6) {Track order};
  \node[uc] (u7) at (0,-2.6) {Manage products and categories};
  \node[uc] (u8) at (0,-3.4) {Process orders / collect COD};
  \node[uc] (u9) at (0,-4.2) {Edit settings and home page};
  \node[uc] (u10) at (0,-5.0) {Manage leads and subscribers};
  \node[uc] (u11) at (0,-5.8) {View analytics and audit log};
  \actor{-5.6,0.9}{cust}{Guest /\\Customer}
  \actor{5.6,-4.2}{admin}{Admin}
  \actor{5.6,0.3}{payu}{PayU\\(external)}
  \foreach \u in {u1,u2,u3,u4,u5,u6} \draw (-5.3,1.2) -- (\u.west);
  \foreach \u in {u7,u8,u9,u10,u11} \draw (5.3,-3.9) -- (\u.east);
  \draw (5.3,0.6) -- (u4.east);
\end{tikzpicture}
\caption{Use case diagram}
\label{fig:uc}
\end{figure}

\subsection{UML Class Diagrams}
\begin{figure}[H]
\centering
\begin{tikzpicture}[font=\scriptsize,
  cls/.style={draw,rectangle split,rectangle split parts=3,fill=black!3,align=left,text width=31mm,inner sep=3pt}]
  \node[cls] (cat) at (-5,3) {\textbf{Category}\nodepart{second}- id, name, slug\\- status\nodepart{third}};
  \node[cls] (prod) at (0,3) {\textbf{Product}\nodepart{second}- id, title, slug\\- status (enum)\nodepart{third}+ fits(model): bool};
  \node[cls] (vm) at (5,3) {\textbf{VehicleModel}\nodepart{second}- id, make, model\nodepart{third}};
  \node[cls] (fit) at (5,0.6) {\textbf{Fitment}\nodepart{second}- productId, modelId\nodepart{third}};
  \node[cls] (sku) at (0,0.4) {\textbf{SKU}\nodepart{second}- id, price, stock\\- dimension values\nodepart{third}+ inStock(): bool};
  \node[cls] (img) at (-5,0.6) {\textbf{ProductImage}\nodepart{second}- url, position\nodepart{third}};
  \node[cls] (user) at (-5,-2.6) {\textbf{User}\nodepart{second}- id, email, phone\\- isAdmin\nodepart{third}+ orders()};
  \node[cls] (ord) at (0,-2.6) {\textbf{Order}\nodepart{second}- orderNumber, status\\- subtotal, discount\\- shippingFee, gst, total\\- paymentMethod\\- advanceAmount\nodepart{third}+ balanceDue()};
  \node[cls] (oi) at (5,-2.6) {\textbf{OrderItem}\nodepart{second}- skuId, title, qty\\- unitPrice (snapshot)\nodepart{third}};
  \node[cls] (ev) at (0,-6.0) {\textbf{OrderEvent}\nodepart{second}- type, note, at\nodepart{third}};
  \draw (cat) -- node[above]{0..* \quad\quad 0..*} (prod);
  \draw (prod) -- node[above]{0..* \quad\quad 0..*} (vm);
  \draw[dashed] ($(prod.east)!0.5!(vm.west)$) -- (fit);
  \draw[-{Diamond[fill=black]}] (sku) -- node[right]{0..* \ 1} (prod);
  \draw[-{Diamond[fill=black]}] (img) -- node[above]{1..* \ 1} (prod.south west);
  \draw (user) -- node[above]{1 \quad\quad 0..*} (ord);
  \draw[-{Diamond[fill=black]}] (oi) -- node[above]{1..* \quad\quad 1} (ord);
  \draw[-{Diamond[fill=black]}] (ev) -- node[right]{0..* \ 1} (ord);
\end{tikzpicture}
\caption{Domain class diagram (selected attributes)}
\label{fig:class}
\end{figure}

\subsection{UML Sequence Diagram}
\begin{figure}[H]
\centering
\begin{tikzpicture}[font=\scriptsize]
  \foreach \x/\n/\t in {0/c/Customer UI,3/a/Hono API,6/d/PostgreSQL,9/p/PayU,12/j/Inngest} {
    \node[box,minimum width=20mm] (\n) at (\x,0) {\t};
    \draw[dashed,black!50] (\x,-0.45) -- (\x,-9.6);
  }
  \draw[arr] (0,-1) -- node[lbl,above=1pt]{1 POST /checkout} (3,-1);
  \draw[arr] (3,-1.7) -- node[lbl,above=1pt]{2 load cart, SKUs, settings} (6,-1.7);
  \draw[arr,dashed] (6,-2.4) -- node[lbl,above=1pt]{3 trusted data} (3,-2.4);
  \draw[arr] (3,-2.9) -- ++(0.5,0) -- ++(0,-0.4) -- ++(-0.5,0); \node[lbl,right] at (3.55,-3.1) {4 compute totals};
  \draw[arr] (3,-3.9) -- node[lbl,above=1pt]{5 create order + items} (6,-3.9);
  \draw[arr,dashed] (3,-4.6) -- node[lbl,above=1pt]{6 signed PayU form} (0,-4.6);
  \draw[arr] (0,-5.3) -- node[lbl,above=1pt,pos=0.35]{7 pay (full or COD advance)} (9,-5.3);
  \draw[arr] (9,-6.0) -- node[lbl,above=1pt,pos=0.5]{8 callback / webhook} (3,-6.0);
  \draw[arr] (3,-6.5) -- ++(0.5,0) -- ++(0,-0.4) -- ++(-0.5,0); \node[lbl,right] at (3.55,-6.7) {9 verify hash + amount};
  \draw (2.2,-7.2) rectangle (12.6,-9.4); \node[anchor=north west,font=\scriptsize\bfseries,draw,fill=white] at (2.2,-7.2) {opt [verified]};
  \draw[arr] (3,-8.1) -- node[lbl,above=1pt]{10 mark paid, add event} (6,-8.1);
  \draw[arr] (3,-8.9) -- node[lbl,above=1pt,pos=0.7]{11 order.paid $\rightarrow$ email + SMS} (12,-8.9);
\end{tikzpicture}
\caption{Checkout and verified-payment sequence}
\label{fig:seq}
\end{figure}

\subsection{UML Activity Diagram}
\begin{figure}[H]
\centering
\begin{tikzpicture}[font=\scriptsize,node distance=5mm]
  \draw (-2,1) rectangle (2,-9.4); \draw (2,1) rectangle (7.6,-9.4); \draw (7.6,1) rectangle (10.6,-9.4);
  \node[font=\footnotesize\bfseries] at (0,0.7) {Customer}; \node[font=\footnotesize\bfseries] at (4.8,0.7) {Application}; \node[font=\footnotesize\bfseries] at (9.1,0.7) {PayU};
  \fill (0,0.2) circle (1mm);
  \node[box] (a1) at (0,-0.6) {Select car and SKU};
  \node[box] (a2) at (0,-1.8) {Sign in and\\enter address};
  \node[box] (a3) at (4.8,-1.8) {Validate cart, pincode};
  \node[box] (a4) at (4.8,-3.0) {Compute subtotal,\\shipping, GST, total};
  \node[box] (a5) at (4.8,-4.3) {Create order (pending)};
  \node[box] (a6) at (9.1,-4.3) {Pay full /\\15\% advance};
  \node[box] (a7) at (4.8,-5.7) {Verify hash and amount};
  \node[draw,diamond,aspect=1.6,inner sep=1pt] (d) at (4.8,-6.9) {valid?};
  \node[box] (n) at (3.2,-8.3) {Keep pending,\\allow retry};
  \node[box] (y) at (6.4,-8.3) {Mark paid,\\notify};
  \draw[arr] (0,0.1) -- (a1); \draw[arr] (a1) -- (a2); \draw[arr] (a2) -- (a3); \draw[arr] (a3) -- (a4);
  \draw[arr] (a4) -- (a5); \draw[arr] (a5) -- (a6); \draw[arr] (a6) |- (a7); \draw[arr] (a7) -- (d);
  \draw[arr] (d) -| node[lbl,pos=0.3]{no} (n); \draw[arr] (d) -| node[lbl,pos=0.3]{yes} (y);
  \draw (n.south) ++(0,-0.25) circle (1.2mm); \fill (n.south) ++(0,-0.25) circle (0.7mm); \draw (n.south) -- ++(0,-0.13);
  \draw (y.south) ++(0,-0.25) circle (1.2mm); \fill (y.south) ++(0,-0.25) circle (0.7mm); \draw (y.south) -- ++(0,-0.13);
\end{tikzpicture}
\caption{Checkout activity diagram with swimlanes}
\label{fig:act}
\end{figure}

\subsection{UML State Machine Diagrams}
Figure~\ref{fig:state} uses the actual \texttt{order\_status} values.
Figure~\ref{fig:claim} is a design-only lifecycle for damage claims
(eligibility: within four days of delivery, with an unboxing video).

\begin{figure}[H]
\centering
\begin{tikzpicture}[font=\scriptsize]
  \fill (-1.7,0) circle (1mm);
  \node[state] (p) at (0,0) {pending};
  \node[state] (pa) at (3.6,0) {paid};
  \node[state] (c) at (7.2,0) {confirmed};
  \node[state] (s) at (10.8,0) {shipped};
  \node[state] (x) at (0,-2.4) {cancelled};
  \node[state] (r) at (5.4,-2.4) {refunded};
  \node[state] (d) at (10.8,-2.4) {delivered};
  \draw[arr] (-1.65,0) -- (p);
  \draw[arr] (p) -- node[lbl,above=2pt]{verified\\payment} (pa);
  \draw[arr] (pa) -- node[lbl,above=2pt]{admin\\confirms} (c);
  \draw[arr] (c) -- node[lbl,above=2pt]{dispatch} (s);
  \draw[arr] (s) -- node[lbl,right=2pt]{received} (d);
  \draw[arr] (p) -- node[lbl,left=2pt]{fail /\\cancel} (x);
  \draw[arr] (pa) -- node[lbl,left=4pt]{refund} (r);
  \draw[arr] (c) -- node[lbl,right=4pt]{refund} (r);
  \draw[arr] ([xshift=-4mm]p.north) to[out=110,in=70,looseness=3] node[lbl,above]{payment failed, retry} ([xshift=4mm]p.north);
  \draw (d.south) ++(0,-0.35) circle (1.2mm); \fill (d.south) ++(0,-0.35) circle (0.7mm); \draw (d.south) -- ++(0,-0.23);
\end{tikzpicture}
\caption{Order state machine (from the \texttt{order\_status} enum)}
\label{fig:state}
\end{figure}

\begin{figure}[H]
\centering
\begin{tikzpicture}[font=\scriptsize]
  \fill (-1.7,0) circle (1mm); \draw[arr] (-1.65,0) -- (-1.15,0);
  \node[state] (r) at (0,0) {RAISED};
  \node[state] (u) at (4.2,0) {UNDER\_REVIEW};
  \node[state] (rj) at (0,-2) {REJECTED};
  \node[state] (ap) at (4.2,-2) {APPROVED};
  \draw[arr] (r) -- node[lbl,above=2pt]{eligible claim} (u);
  \draw[arr] (r) -- node[lbl,left=2pt]{ineligible} (rj);
  \draw[arr] (u) -- node[lbl,above=2pt]{evidence insufficient} (rj);
  \draw[arr] (u) -- node[lbl,right=2pt]{evidence sufficient} (ap);
\end{tikzpicture}
\caption{Damage-claim state machine (design only)}
\label{fig:claim}
\end{figure}

\subsection{Data Model}
At revision \texttt{fb14d3d} the schema had 16 migrations. The final schema
(28 tables) is listed in the final report.

\chapter{Results and Evaluation}

\section{Testing Methodology}

\begin{itemize}
  \item \textbf{Unit testing:} Vitest suites for pricing, hashing, cart state,
        checkout helpers, category logic and SKU-matrix calculations.
  \item \textbf{Integration testing:} API route tests, API-client tests, MSW
        handlers and an end-to-end checkout smoke script.
  \item \textbf{Security testing:} an OWASP-mapped test catalogue covering
        ownership, JWTs, payment integrity, injection, uploads and business
        logic.
  \item \textbf{Performance testing:} a search benchmark script; live
        measurements were deferred to the final evaluation.
  \item \textbf{User testing:} not performed at this stage.
\end{itemize}

\section{Performance Metrics}

\begin{table}[H]
\centering\small
\begin{tabularx}{\linewidth}{l l X}
\toprule
\textbf{Metric} & \textbf{Value (15 Oct 2026)} & \textbf{Note} \\
\midrule
Commits on main & 75 & Since 31 May 2026. \\
Test files & 66 (44 API, 22 storefront) & Files present at that revision. \\
Database migrations & 16 & \\
Latency / throughput & not measured & Measured for the final report. \\
User evaluation & not performed & \\
\bottomrule
\end{tabularx}
\caption{Project scale at the prototype stage}
\end{table}

\section{Analysis}

The prototype met the functional objectives: vehicle-aware discovery,
customer and admin access, server-owned checkout and administrative
operations. Two deployment problems shaped the design. Refresh cookies set by
the API domain were blocked in browsers that block third-party cookies, so
session routes were moved behind a same-origin proxy. The Vercel image
optimisation quota ran out, so image resizing moved to Supabase's rendering
endpoint. The remaining work at this stage was to measure performance, run
the suites against a real database and PayU sandbox, and decide whether
claims and courier automation belonged in scope.

\chapter{Conclusion}

\section{Summary of Work}
By mid-semester the team had a deployed vehicle-accessory store with a
Next.js storefront, a Hono API, PostgreSQL persistence and PayU payments,
covering fitment, variants, authentication, checkout, order operations,
administration, analytics and configurable branding.

\section{Key Findings}
\begin{itemize}
  \item \textbf{Finding 1:} Fitment belongs to the product--vehicle
        relationship, and the SKU is what an order line buys.
  \item \textbf{Finding 2:} Checkout must use trusted server-side data, and
        payment confirmation must depend on verified gateway evidence.
  \item \textbf{Finding 3:} Hosting limits (cookies, image quotas) changed the
        architecture as much as the requirements did.
\end{itemize}

\section{Future Work and Recommendations}
\begin{enumerate}
  \item Guest checkout and promo codes (delivered after mid-semester).
  \item Measure page speed and API latency on the live site.
  \item Decide on the damage-claim workflow and courier integration.
  \item Verify payment replay handling and stock contention under load.
\end{enumerate}

\section{Final Remarks}
The prototype established the main customer and staff workflows; the next
milestone was to attach measured evidence to them, which the final report
does.

\begin{thebibliography}{SRS26}
\bibitem[SRS26]{srs} V.~Verma, R.~Verma, P.~Sharma, T.~Mohindru. \emph{Software Requirements Specification for SV Car Customs}. UCS503, TIET, 2026.
\bibitem[Next25]{next} Vercel. \emph{Next.js Documentation}. \url{https://nextjs.org/docs}, 2025.
\bibitem[Hono25]{hono} Hono. \emph{Hono Documentation}. \url{https://hono.dev}, 2025.
\bibitem[PayU25]{payu} PayU India. \emph{PayU Hosted Checkout Integration}. \url{https://docs.payu.in}, 2025.
\bibitem[OWASP21]{owasp} OWASP Foundation. \emph{OWASP Top 10: 2021}. 2021.
\end{thebibliography}

\end{document}
